\chapter{Evaluation Harnesses}
\label{chapter:Evaluation}

A research agent that cannot be evaluated cannot be trusted.\index{Evaluation harness}
Evaluation is not a post-hoc check performed after a system is built; it is the foundation on which every comparison, every improvement claim, and eventually every evolutionary search rests.
This chapter develops the theory and practice of the evaluation harness — the reusable module that scores an agent's output against ground truth.

The harness introduced in this chapter is the same object that becomes the fitness function of the evolutionary framework in Chapter~\ref{chapter:EvolutionaryFrameworks}.
Build it cleanly; you will live with it for the rest of the course.

\section{The Evaluation Principle}
\label{section:EvaluationPrinciple}

An \emph{evaluation harness}\index{Evaluation harness} is a function that takes a candidate artifact and a reference dataset and returns a numerical score.
\begin{definition}
\label{definition:EvalHarness}
An \emph{evaluation harness} is a tuple $(\mathcal{D}, \Score)$ where
\begin{itemize}
\item $\mathcal{D} = (\mathcal{D}_{\mathrm{train}}, \mathcal{D}_{\mathrm{test}})$ is a partition of the data into a training split and a held-out test split, and
\item $\Score : \Artifacts \times \mathcal{D}_{\mathrm{test}} \to \RealNumbers$ is a scoring function that evaluates a candidate artifact $a \in \Artifacts$ against the test split.
\end{itemize}
The score of artifact $a$ is $\Score(a, \mathcal{D}_{\mathrm{test}})$.
\end{definition}

The separation of training and test data, and the requirement that the score be computed on the test split alone, are the most important properties of the harness.
They are also the easiest to violate accidentally.

\subsection{Why Evaluation Must Be Automatic}

Automatic evaluation — a numerical score computed without human judgment — has three decisive advantages for agentic and evolutionary systems.
\begin{enumerate}
\item \textbf{Speed.} An automated harness scores hundreds of candidates in the time a human would score one.
\item \textbf{Consistency.} Automated scores are deterministic; human judgments vary.
\item \textbf{Composability.} A numerical score can be directly compared, optimized, and used as a fitness signal.
\end{enumerate}
These properties are requirements for the evolutionary half of the course.
An evaluation that requires human intervention cannot drive an automated search.

\subsection{What Counts as a Good Score}

A good score is:
\begin{itemize}
\item \textbf{Faithful}: high scores correspond to genuinely good artifacts, not to artifacts that exploit loopholes in the scoring function.
\item \textbf{Discriminating}: the score distinguishes between candidates of different quality; all candidates scoring the same is as bad as no score at all.
\item \textbf{Leakage-resistant}: the score cannot be improved by exposing the artifact to test data during training.
\item \textbf{Stable}: small, inconsequential changes to the artifact produce small changes in the score; large changes in quality produce large changes in the score.
\end{itemize}

\section{Held-Out Data and Leakage}
\label{section:Leakage}

\subsection{The Train/Test Split}

The fundamental practice of supervised evaluation is to partition the data before any modeling begins and to hold the test split entirely in reserve.\index{Train/test split}

\begin{principle}
\label{principle:TestSplitInviolable}
The test split is inviolable.
It must not be used to select hyperparameters, tune the model, or make any decision about the artifact.
The moment any information about the test split influences the artifact, the test score is optimistic and the comparison is compromised.
\end{principle}

In practice: load the full dataset, split immediately (before looking at the data), and pass only $\mathcal{D}_{\mathrm{train}}$ to the agent.
The test split is sealed inside the harness; the agent never sees it.

\subsection{What Leakage Is}

\emph{Leakage}\index{Leakage} is the unintended exposure of test-set information to the artifact during its construction.
Leakage inflates test scores, making a weak artifact appear strong.
Its consequences are especially severe when a harness is used as a fitness function: an evolutionary search that optimizes a leaking harness will produce artifacts that score well on leaked tests and poorly in the real world.

\begin{example}
Common forms of leakage, in increasing subtlety:
\begin{enumerate}
\item \textbf{Scoring on training data.} The function returns error on $\mathcal{D}_{\mathrm{train}}$ rather than $\mathcal{D}_{\mathrm{test}}$.
An overfit artifact achieves near-zero training error and receives an excellent score.
\item \textbf{Feature construction from the full dataset.} Computing a normalization mean from the full dataset before splitting contaminates the training split with information from the test split.
\item \textbf{Adaptive re-splitting.} Trying multiple splits and selecting the one with the best score on a specific candidate.
The resulting split is no longer independent of the artifact.
\end{enumerate}
\end{example}

\subsection{Leakage-Resistant Design}

A leakage-resistant harness separates the split from the fit from the score into three independent stages:
\begin{lstlisting}[language=python]
def build_harness(data, test_fraction=0.2, seed=42):
    """Split data and return a scorer sealed over the test set."""
    X, y = data['X'], data['y']
    idx = np.random.default_rng(seed).permutation(len(y))
    n_test = int(len(y) * test_fraction)
    train_idx, test_idx = idx[n_test:], idx[:n_test]
    X_test, y_test = X[test_idx], y[test_idx]

    def score(model):
        """Score model on the sealed test set."""
        y_pred = model.predict(X_test)
        return float(np.mean((y_pred - y_test) ** 2))  # MSE

    return score   # the scorer is a closure over X_test, y_test
\end{lstlisting}
The caller receives only a scoring function; the test data itself is inaccessible.
This is the design pattern used in the course harness.

\section{Scoring Functions}
\label{section:ScoringFunctions}

\subsection{Predictive Error}

The most natural score for a model discovery task is \emph{predictive error}: how well does the discovered model predict the test targets?
For a model $\hat{m}$ and test data $\{(x_i, y_i)\}_{i=1}^{n}$, the \emph{mean squared error} (MSE) is
\begin{equation}
\mathrm{MSE}(\hat{m}) = \frac{1}{n} \sum_{i=1}^{n} \bigl( \hat{m}(x_i) - y_i \bigr)^2 .
\label{equation:MSE}
\end{equation}
MSE penalizes large errors quadratically and is differentiable, which makes it attractive for gradient-based methods.
When MSE is on the leaderboard, the researcher reports the lower value as the better one.

\subsection{Parsimony}

A model that perfectly predicts the test data by memorizing it is not a discovery.
\emph{Parsimony}\index{Parsimony} — the preference for simpler models among those of comparable predictive quality — guards against this failure mode.

The parsimony of a symbolic expression is typically measured by its \emph{complexity}: the number of nodes in its expression tree, or the number of operations required to evaluate it.
A simpler expression is easier to interpret, more likely to generalize, and less likely to be an accidental fit.

\begin{example}
The expressions $f_1(x) = x^2$ and $f_2(x) = x^2 + 10^{-9} x^7$ may achieve similar MSE on a clean test set, but $f_1$ has complexity~2 and $f_2$ has complexity~5.
Parsimony selects $f_1$ as the discovery.
\end{example}

\subsection{Multi-Objective Scoring}
\label{subsection:MultiObjectiveScoring}

A composite score combines accuracy and parsimony into a single number.
A common form is
\begin{equation}
\Score(\hat{m}) = \mathrm{MSE}(\hat{m}) + \lambda \cdot C(\hat{m}) ,
\label{equation:CompositeScore}
\end{equation}
where $C(\hat{m})$ is the complexity of $\hat{m}$ and $\lambda > 0$ is a parsimony weight.
Lower scores are better.
The weight $\lambda$ encodes the tradeoff between accuracy and simplicity; the course benchmarks specify its value so that all participants use the same harness.

\begin{remark}
The parsimony weight $\lambda$ is a design choice, not a universal constant.
Too small a $\lambda$ and the optimizer ignores complexity; too large and it ignores accuracy.
The appropriate value depends on the scale of MSE (which changes with the dataset) and the expected complexity of the target expression.
For the course benchmarks, $\lambda$ is calibrated so that a model twice as complex must reduce MSE by a meaningful amount to be preferred.
\end{remark}

\section{Evaluation Hacking}
\label{section:EvalHacking}

\emph{Evaluation hacking}\index{Evaluation hacking} (also known as \emph{Goodhart's Law}\index{Goodhart's Law} in the context of learned systems) is the phenomenon in which a system optimizes the \emph{proxy} (the score) rather than the \emph{goal} (genuine discovery of a correct model).
Wherever a score is computed automatically and an optimizer is pointed at it, evaluation hacking will occur if there is any gap between the score and the goal.

\begin{principle}
\label{principle:GoodhartsLaw}
Any metric that becomes a target ceases to be a good metric.
The stronger and more persistent the optimizer, the more aggressively it will find and exploit weaknesses in the scoring function.
Designing a harness you would trust an adversary to attack is more reliable than designing one that looks good on the first round of benign candidates.
\end{principle}

Common attack surfaces in a model discovery harness:
\begin{enumerate}
\item \textbf{Trivial re-parameterization.} A parsimony term that counts nodes rewards replacing $2 \cdot x$ with a single constant $c$ learned from the data (zero nodes for $c$, but two nodes for $2 \cdot x$), even though both represent the same model.
\item \textbf{High-degree polynomial fitting.} A sufficiently high-degree polynomial can achieve near-zero MSE on any finite dataset — including the test set, if the harness leaks.
\item \textbf{Memorizing the test set via complexity.} An artifact that memorizes $n$ test points as a piecewise constant function with $n$ pieces has zero MSE and complexity $O(n)$ — which may still produce a low composite score if $\lambda$ is too small.
\end{enumerate}

Recognizing these attack surfaces in advance and closing them in the harness design is the primary defense.

\section{The Harness as a Contract}
\label{section:HarnessContract}

The evaluation harness is used across Challenges~2, 3, and 4, and becomes the fitness function in the capstone.
Its interface must be \emph{stable} across this span.

\begin{definition}
\label{definition:HarnessContract}
A harness interface is a pair $(\texttt{score}, \texttt{report})$ where
\begin{itemize}
\item $\texttt{score}(\text{model}) \to \text{float}$ returns the composite score (lower is better), and
\item $\texttt{report}(\text{model}) \to \text{dict}$ returns a structured breakdown: MSE, complexity, and score.
\end{itemize}
A harness satisfying this interface can be dropped into any unit in the course without modification.
\end{definition}

If the harness interface must change — for example, because a new scoring dimension is added in a later challenge — the change must be documented and applied consistently to all prior baselines.
An inconsistent harness produces incomparable leaderboard entries.

\section*{Further Reading}

\begin{small}
\begin{enumerate}
\item \texttt{benchmarks/README.md} in the course repository — the specification of the shared harness, including the parsimony weight and the leaderboard format.
\item Kapoor, S., and Narayanan, A., ``Leakage and the Reproducibility Crisis in Machine Learning,'' \emph{PLOS ONE}, 2023 — an empirical survey of leakage in published ML results.
\item Sculley, D., et al., ``Challenging Common Assumptions in the Unsupervised Pretraining of Deep Neural Networks,'' \emph{ICML}, 2014 — a classic treatment of evaluation pitfalls in ML research.
\end{enumerate}
\end{small}
